\section{Discussion - \textit{Ultra as a Deployed Application and Its Generality}}
In this section, we discuss the characteristics of ULTRA as required for a deployment track paper but not discussed elsewhere. We describe its development experience leading to deployment at a US university. We also demonstrate its generality by considering an altogether different setting from India but leading to similar results: with increased sophistication of methods, the quality of teams increase and the amount of recommended teams decrease. 

\subsection{Development and Deployment}
The development of ULTRA first began in Summer 2021 with a team of 8 developers and an initial prototype was created within 3 months with M0 and M1 for user feedback (August 2021). It included a very small dataset of researchers and RFPs, a single method (i.e., a string-based matching algorithm and greedy teaming strategy), no goodness score and instead only a match percentage threshold, and a smaller-scale empirical evaluation with a few experts at a single University that indicated the promise of such a tool. Once the pilot study results seemed promising, ULTRA was refreshed with additional data; re-imagined and re-designed in Fall 2022; improved through iterative review, feedback, and testing; enhanced with \textit{M2} and \textit{M3,} and alpha-tested with users from different departments at the University of South Carolina for a year. The system was deployed on May 1, 2023 and evaluated under an IRB-approved protocol for 28 days, and the results are as reported in the paper. It remains publicly available as of writing this paper (November 2023). %for use till date. 
During July-August 2023, we evaluated ULTRA with data from another setting in a different country: researchers at Indian Institute of Technology-Roorkee (IIT-R), India and calls from India's funding agencies.   

One main challenge during development was access to clean data related to RFPs and relevant researcher profiles. Due to inconsistencies in data formatting, missing values, and irrelevant or out-of-date entries, exploring automated approaches had been unsuccessful, and data cleaning had  only been possible through iterative collaborations. 
%As such, the datasets used were also smaller in number. 
After deployment, additional challenges were raised. One challenge was maintaining data integrity across changes to ULTRA's servers and infrastructure, and another was change in researchers over time due to hiring or attrition. It was essential to run frequent tests to enhance user experience, create a working feedback loop, and ensure a scalable architecture with minimal overhead. 

\begin{table}
    \centering
    \begin{tabular}{c|c|c}
    \toprule 
    Method  & Average Quality & Average Volume \\
    \midrule 
    \textit{M0} & $0.0896\pm0.0006$ & $10$ \\
    \textit{M1} & $0.4218\pm0.0011$ & $8$ \\
    \textit{M2} & $0.4292\pm0.0017$ & $7$ \\
    \textit{M3} & $0.5835\pm0.0203$ & $1$ \\
    \bottomrule
    \end{tabular}
    \caption{Average quality ($G$) and volume of teams (\#$T$) shown \textit{per} researcher ($r_j$) at IIT-R. This was done for each method $M_i$, across 100 RFPs and 46 researchers. For average quality, we report the mean and standard deviation, denoted as \texttt{mean$\pm$STD}.} \label{tab:volume_quality_iitr}
\end{table}

\subsection{Generalizing to Second Institution: IIT-R}
We additionally extended ULTRA to another university in a different region of the world: Indian Institute of Technology-Roorkee (IIT-R), India. We gathered publicly available data on RFPs from the Department of Science \& Technology (DST), a division of Research and Development (R\&D) programmes belonging to and funded by the Government of India's Ministry of Science and Technology~\cite{dst2023department}. We further collected data about IIT-R's faculty members and their respective research profiles. 
%Due to the challenges described above, 
Our initial evaluation is with 100 RFPs and 46 researchers.
Table~\ref{tab:volume_quality_iitr} shows the computational evaluation of ULTRA on IIT-R's data. 
% Due to additional time constraints, we were not able to perform a human study with a larger dataset.
From the quantitative results, we observe a similar trend to Table~\ref{tab:volume_quality}, where teams formed algorithmically led to an increased precision and quality, and visible decrease in quantity.
Since human study at any institution is subject to local policies and workplace culture, we leave performance of such a study at IIT-R as a future work.

